(1) 两无磁、各向同性、无损介质在 \(z=0\) 交界：上方 \(z>0\) 为实正介电常数 \(\varepsilon_d\)，下方 \(z<0\) 为实负 \(\varepsilon_m\)，均表示相对介电常数。取时空因子 \(e^{i\beta x-i\omega t}\)，考虑沿界面传播、离界面指数衰减的TM模式，令 \(k_0=\omega/c\)。写出两侧场及衰减常数 \(\kappa_j\) 的定义。
(2) 利用 \(E_x,H_y\) 的切向连续性推导表面模式本征方程，求 \(\beta(\omega)\)，明确实束缚模式存在的参数条件。
(3) 求两侧场的穿透长度 \(1/\kappa_d,1/\kappa_m\)。若 \(\varepsilon_d=1\) 且金属满足无损Drude \(\varepsilon_m(\omega)=1-\omega_p^2/\omega^2\)，求大 \(\beta\) 极限的表面等离激元频率和低频光线渐近。
(4) 说明高 \(\beta\) 极限群速度趋势及金属微弱吸收时传播常数的变化，讨论色散介质电磁能密度为何不能直接用常数 \(\varepsilon_m|E|^2/4\) 计算。
